% ============================================================
% AUTOMATISCH GEGENEREERD BESTAND
% Niet handmatig aanpassen.
% ============================================================

\ifdefined\HCode
  % Geen afzonderlijke module-openingspagina in HTML.
\else
  \FVDmoduleopening
    {De digitale wereld begrijpen}
    {In deze module bouwen we de wiskundige basis op die je nodig hebt voor de rest van de cursus. We vertrekken van getallen en rekenvaardigheid en groeien stap voor stap naar algebra, formules en eenvoudige wiskundige modellen. \\ \ \\ We oefenen hoofdrekenen, negatieve getallen, kommagetallen, breuken, procenten en verhoudingen. Daarna onderzoeken we machten, wortels en wetenschappelijke notatie. Vanuit concrete berekeningen maken we vervolgens de stap naar rekenen met letters en formules. \\ \ \\ De wiskunde krijgt betekenis in toepassingen met onder andere pixels, resoluties, beeldverhoudingen, coördinaten, schaalfactoren en andere grootheden uit digitale creatie en technologie. Waar het zinvol is, gebruiken we GeoGebra om verbanden te onderzoeken en zichtbaar te maken. \\ \ \\ Een computer kan razendsnel rekenen, maar jij moet kunnen beoordelen of een resultaat logisch is. Daarom leren we ook schatten, controleren, een geschikte strategie kiezen en kritisch naar resultaten kijken. \begin{center} \textbf{Het doel is niet alleen correct rekenen, maar wiskunde gebruiken om verbanden te begrijpen, problemen op te lossen en de digitale wereld te beschrijven.} \end{center}}
    {%
    \FVDmodulethemetile{1}{Rekenen in een digitale wereld}
    \FVDmodulethemetile{2}{Machten, wortels en grootteorde}
    \FVDmodulethemetile{3}{Algebra als taal}
    \FVDmodulethemetile{4}{Wiskunde onderzoeken en toepassen}
    }
\fi
